% ============================================================
% IEMOCAP 50% Limitation — Discussion Section
% ============================================================

\paragraph{Limitation Analysis: High-Label Regime on Small Datasets.}
We observe that ECR's advantage diminishes in the high-label regime (50\% labeled data) 
on smaller datasets such as IEMOCAP. Specifically, at 50\% label ratio, FixMatch achieves 
$59.65 \pm 0.79\%$ WF1 compared to ECR's $59.12 \pm 0.56\%$ ($p = 0.018$). We attribute 
this to two factors:

\begin{enumerate}
    \item \textbf{Reduced unlabeled pool.} At 50\% label ratio on IEMOCAP (only 120 training 
    dialogues), the unlabeled pool shrinks to $\sim$60 dialogues — insufficient for the 
    KL-based consistency regularization to provide meaningful gradients.
    
    \item \textbf{Diminishing returns of uncertainty weighting.} When labeled data is abundant, 
    the supervised evidential loss already produces well-calibrated Dirichlet parameters, 
    leaving less room for ECR's certainty-weighted consistency to contribute.
\end{enumerate}

Notably, ECR's primary design target — \textit{low-label federated settings} — shows 
consistent and statistically significant improvements: $+7.47\%$ at 5\% labels 
(IEMOCAP) and $+0.48$--$0.60\%$ across all MELD conditions ($p < 0.02$). This aligns 
with ECR's theoretical motivation: when pseudo-labels are unreliable (low-label regime), 
KL-based consistency avoids the confirmation bias inherent in hard-threshold methods 
like FixMatch.

We further verified through a grid search over $\lambda_u \in \{0.3, 0.5, 1.0, 2.0, 3.0, 5.0\}$ 
and warm-up schedules $\{5, 10, 20\}$ rounds (54 configurations) that this gap cannot be 
closed by hyperparameter tuning alone, confirming it as a structural characteristic of 
consistency-based SSL in data-rich regimes rather than a tuning artifact.
